\documentclass[aspectratio=43]{beamer}
\usetheme{Boadilla}
\setbeamertemplate{navigation symbols}{}
\usepackage{tikz}
\usepackage{amsmath}
\usepackage{stmaryrd}
\usepackage{mathtools}

\title[$\mu$-Calculus 101]{The Madness of the Modal $\mu$-Calculus}
\subtitle{Not your granny's syntax with binding!}
\author[S. Watters]{Sean Watters}
\institute{University of Strathclyde}
\date{28th Oct 2024}


\newcommand{\meaningof}[1]{\ensuremath{\llbracket #1 \rrbracket_{\mathcal{M}}}}

\begin{document}


\begin{frame}
  \titlepage{}
\end{frame}

\begin{frame}{Plan}
\begin{enumerate}
\item Introduction to the $\mu$-calculus.
        \smallskip
\begin{enumerate}
  \item Model Checking
        \smallskip
  \item Kripke Semantics
        \smallskip
  \item Strict Positivity
        \smallskip
  \item The Closure (enter the madness)
\end{enumerate}
        \smallskip
\item Introduction to well-scoped De Bruijn syntax.
        \smallskip
\begin{enumerate}
  \item Thinnings
        \smallskip
  \item Weakening
        \smallskip
  \item Parallel substitution
\end{enumerate}
        \smallskip
\item A peek at \emph{sublimely-scoped} De Bruijn syntax (if time permits).
\end{enumerate}
\end{frame}

\begin{frame}{Logic as a Specification Language (1)}
  A \textbf{logic} is a formal system for making statements of fact.
  eg: propositional logic, first-order logic, etc.

  \bigskip{}

  Traditionally, we define the meaning of formulae relative to some mathematical structure called a \textbf{model}.
  (Model-theoretic semantics).

  \bigskip{}

  The \textbf{model checking problem} for a logic is the problem of deciding whether, given a formula and model of the logic, the formula is validated in that model.

  \bigskip{}

  Model Checking as a \textbf{formal verification} technique involves representing the system undergoing verification as a model, and the properties being verified as formulas.
  As long as the model checking problem for the logic is decidable, we can algorithmically verify that the property is true.
\end{frame}

\begin{frame}{Logic as a Specification Language (2)}

  Two questions:
  \bigskip{}
  \begin{enumerate}
    \item What is the right notion of model for real-world systems?
    \smallskip
    \item Which logics (best) let us reason about such models?
  \end{enumerate}

\end{frame}

\begin{frame}{Propositional Logic}
  Given some set $A$ of propositional atoms, for all $a \in A$, let PL be the set of terms generated by:
  \begin{equation*}
    \varphi := a~|~\neg \varphi~|~\varphi \land \varphi~|~\varphi \lor \varphi~|~\varphi \Rightarrow \varphi
  \end{equation*}

  Models are valuation functions $A \to 2$.

  \bigskip{}

  \begin{block}{Propositional Logic for Model Checking?}
  \begin{itemize}
          \smallskip
    \item \textbf{The Good}: Model checking problem is decidable.
          \smallskip{}
    \item \textbf{The Bad}: Not very expressive; purely propositional encodings are clunky.
  \end{itemize}
  \end{block}

  \bigskip

  How about first-order logic? Much more expressive, but too expressive! Not decidable!
\end{frame}

\begin{frame}{Kripke Semantics (1)}

  Key idea: Computers and programs tend to progress in discreet time steps, with their properties potentially changing at each step.

  Let's consider models that reflect that!

  % TODO - pictures of LTL infinite stream model, CTL infinitely descending tree model, and a DFA/graph kripke model

\end{frame}

\begin{frame}{Kripke Semantics (2)}

\begin{definition}
  A \textbf{Kripke model} is a tuple $(S,T,V)$ of a set of states $S$, a transition relation $T : S \to S \to \mathsf{Prop}$, and a valuation function $V : S \to A \to \mathsf{Prop}$ (given some set $A$ of atomic propositions).
\end{definition}

\bigskip

Notice that:
\begin{itemize}
  \item $S$ and $T$ induce a (possibly infinite) graph-like structure.
        \smallskip
  \item Each state $s : S$ is equipped with a model of propositional logic, $V~s$.
        \smallskip
  \item Partial application of $T$ on some state $s : S$ yields the ``is a successor of $s$'' predicate.
\end{itemize}

\end{frame}

\begin{frame}{Modal Logic (1)}
  For all propositional atoms $a \in A$, let ML be the set of terms generated by:
  \begin{equation*}
    \varphi := a~|~\neg \varphi~|~\varphi \land \varphi~|~\varphi \lor \varphi~|~\varphi \Rightarrow \varphi~|~\Box \varphi~|~\Diamond \varphi
  \end{equation*}

  \bigskip

  Given a Kripke model $\mathcal{M} = (S,T,V)$, let $\meaningof{-}~- : ML \to S \to \mathsf{Prop}$, where:
  \begin{align*}
    \meaningof{a}~s &:= V~s~a \\
    \meaningof{\lnot \varphi}~s &:= \lnot~(\meaningof{\varphi}~s) \\
    \meaningof{\varphi \land \psi}~s &:= (\meaningof{\varphi}~s) \times (\meaningof{\psi}~s) \\
    \meaningof{\varphi \lor \psi}~s &:= (\meaningof{\varphi}~s) + (\meaningof{\psi}~s) \\
    \meaningof{\varphi \Rightarrow \psi}~s &:= (\meaningof{\varphi}~s) \to (\meaningof{\psi}~s) \\
    \meaningof{\Box \varphi}~s &:= \forall t \in (T~s).~\meaningof{\varphi}~t \\
    \meaningof{\Diamond \varphi}~s &:= \exists t \in (T~s).~\meaningof{\varphi}~t
  \end{align*}
\end{frame}

\begin{frame}{Modal Logic (2)}

\begin{block}{ML for Model Checking?}
\begin{itemize}
  \item \textbf{Good}: Model checking problem is still decidable.
        \smallskip{}
  \item \textbf{Good}: Kripke models are well-suited to representing real systems.
        \smallskip{}
  \item \textbf{Bad}: Expressivity is still limited; can only reason about fixed, finite time steps.
\end{itemize}
\end{block}

\bigskip

We cannot express, for example:

\begin{itemize}
  \item $\varphi$ is always true.
  \item $\varphi$ eventually becomes true, in at least one possible future.
  \item $\varphi$ remains true at least until $\psi$ becomes true, in all possible futures.
\end{itemize}

\bigskip

Reasoning about infinte or unbounded behaviour is problematic.

\bigskip

There's a neat way around that\ldots

  % more expressive logics for kripke models: LTL, CTL, PDL
\end{frame}

\begin{frame}{The Modal $\mu$-Calculus}
  For all propositional atoms $a \in A$ and variable names $x$, let $\mu$ML be the set of terms generated by:
  \begin{equation*}
    \varphi := a~|~\neg a~|~\varphi \land \varphi~|~\varphi \lor \varphi~|~\Box \varphi~|~\Diamond \varphi~|~\mu x. \varphi~|~\nu x. \varphi
  \end{equation*}

  \begin{align*}
    \meaningof{\mu x.~\varphi}~s &:= \mu~\meaningof{\varphi}~t \\
    \meaningof{\Diamond \varphi}~s &:= \exists t \in (T~s).~\meaningof{\varphi}~t
  \end{align*}
  % subsumes all of the above.
  % note: strict positivity
  % mention semantics; classical, game, container
\end{frame}

\begin{frame}{Fixpoint Unfolding}
  % it's the whole point
  % substitution (hand wavy, will do properly with dB in part 2)
\end{frame}

\begin{frame}{The Closure}
  % direct definition...terminating?
  % insability under alpha equivalence
  % relationship to model checking
\end{frame}

% expansion map, and closure via that? no, probably should skip that
%
% \begin{frame}{De Bruijn Syntax}
%   % basic, non-dependent version
% \end{frame}
%
% \begin{frame}{Well-Scoped De Bruijn Syntax}
%   % now with fams and fins
% \end{frame}
%
% \begin{frame}{Parallel Substitution}
%   % the standard formulation of substitution
%   % definition of unfolding
% \end{frame}
%
% \begin{frame}{Thinnings}
%   % define thinnings
%   % define embed
% \end{frame}
%
% \begin{frame}{Weakening}
% \end{frame}
%
% \begin{frame}{The Challenge}
%   % trying to prove correctness of the closure is hard;
%   % trying to reason about a graph-like thing in a tree-like manner is a bad smell;
%   % would need all sorts of painful ad-hoc bookkeeping.
%   % key insight: lets bake the data into the type instead.
%   % also mention: rational data, Hamana's inductive graphs as trees with backedges and sharing
% \end{frame}
%
% \begin{frame}{Fancier Contexts}
%   % the fancier contexts
% \end{frame}
%
% \begin{frame}{Sublimely-Scoped De Bruijn Syntax}
%   % our fancier formulas
% \end{frame}


%
%\begin{frame}{Kripke Semantics}
  %Propositional logic:
  %\begin{itemize}
    %\item Has good decidability/computational properties.
    %\item Is not very expressive.
  %\end{itemize}
%
  %\bigskip
%
  %We can extend it to FO logic, but lose decidability.
%
  %\bigskip
%
  %Idea: What if we could extend it in other ways, while keeping
%
  %Kripke Semantics:
  %\begin{itemize}
    %\item Given a set $P$ of propositions, a model of propositional logic is a function $P \to 2$.
  %\end{itemize}
%
  %\begin{tikzpicture}[->,shorten >=1pt,auto,node distance=2cm,
                    %semithick]
  %\tikzstyle{every state}=[fill=red,draw=none,text=white]
%
  %\node[initial,state] (A)                    {$Q_{0}$};
  %\node[state]         (B) [above right of=A] {$Q_{1}$};
  %\node[state]         (D) [below right of=A] {$Q_{2}$};
  %\node[state]         (C) [below right of=B] {$Q_{3}$};
  %\node[state]         (E) [below of=D]       {$Q_{4}$};
%
  %\path (A) edge              node {} (B)
            %edge              node {} (C)
        %(B) edge [loop above] node {} (B)
            %edge              node {} (C)
        %(C) edge              node {} (D)
            %edge [bend left]  node {} (E)
        %(D) edge [loop below] node {} (D)
            %edge              node {} (A)
        %(E) edge [bend left]  node {} (A);
  %\end{tikzpicture}
%\end{frame}
%
%\begin{frame}{The Propositional Modal $\mu$-calculus}
  %Given a set of variable names $P$ with $p,x \in P$ :
%\begin{equation*}
%L_{\mu}~:=~\top~|~\bot~|~p~|~\lnot p~|~\varphi \land \psi~|~\varphi \lor \psi~|~\Box \varphi~|~\Diamond \varphi~|~\mu{}x.\varphi{}~|~\nu{}x.\varphi{}
%\end{equation*}
%
%\bigskip
%
%Essentially, this is:
%\begin{itemize}
  %\item Propositional logic,
  %\item Plus the modal operators $\Box$ and $\Diamond$,
  %\item Plus the least and greatest fixpoint operators $\mu$ and $\nu$.
%\end{itemize}
%\end{frame}
%
%\begin{frame}{Fixpoint Unfolding}
  %For any fixpoint formula of the form $\mu{}x.\varphi{}$, we have:
  %\begin{equation*}
    %\mu{}x.\varphi \equiv \varphi[x~/~\mu{}x.\varphi{}]
  %\end{equation*}
  %(And dually for $\nu$.)
%
  %\bigskip
%
  %eg:
  %\begin{align*}
    %&\nu{}x.p \land \Box{}x \\
    %\equiv~&p \land \Box{}(\nu{}x.p \land \Box{}x) \\
    %\equiv~&p \land \Box{}(p \land \Box{}(\nu{}x.p \land \Box{}x)) \\
    %\equiv~&\ldots{}
  %\end{align*}
%
  %\bigskip
%
  %Surprising(?) fact: For any formula $\varphi$, take the set of formulas closed under direct subformulas of propositional and modal formulas, and unfoldings of fixpoints. This set is always finite.
%\end{frame}
%
%\begin{frame}{The Fischer-Ladner Closure}
  %\begin{definition}[The Closure of $\phi$]
    %The least set of formulas which contains $\phi$ and is closed under taking direct subformulas of propositional and modal subformulas, and unfolding fixpoint subformulas.
  %\end{definition}
  %\bigskip
%
  %The closure is highly important in the study of the $\mu$-calculus, but\ldots{}
%
  %\begin{center}
  %\textbf{It is not invarient under $\alpha$-equivalence!}
  %\end{center}
%
  %\bigskip
%
  %For example, we have:
  %\begin{align*}
  %&\mu{}x.(\Diamond{}x \land \Box{}(\nu{}y.\mu{}z.(\Diamond{}z \land \Box{}y)))\\
  %\equiv~&\mu{}x.(\Diamond{}x \land \Box{}(\nu{}y.\mu{}x.(\Diamond{}x \land \Box{}y)))
  %\end{align*}
%
  %But their closures are different!
%
%\end{frame}
%
%\begin{frame}{De Bruijn to the Rescue?}
  %\snippetwellscoped{}
  %\snippetscope{}
%\end{frame}
%
%\begin{frame}{\sout{Well} \emph{Sublimely}-Scoped Formulas}
  %\snippetsublimelyscoped{}
%\end{frame}

\end{document}
